\documentclass[11pt,a4paper]{article}

\usepackage[utf8]{inputenc}
\usepackage[T1]{fontenc}
\usepackage[french]{babel}
\usepackage{geometry}
\usepackage{graphicx}
\usepackage{hyperref}
\usepackage{titlesec}
\usepackage{enumitem}

% Configuration simple
\geometry{left=2.5cm, right=2.5cm, top=3cm, bottom=3cm}
\hypersetup{colorlinks=true, linkcolor=blue, urlcolor=blue}

\title{\textbf{Credit Platform}\\Rapport Technique}
\author{Équipe de Développement}
\date{\today}

\begin{document}

\maketitle
\tableofcontents
\newpage

\section{Introduction}

\subsection{Problématique}
Les institutions financières font face à un défi majeur dans la gestion des demandes de crédit provenant de multiples canaux (Email, WhatsApp, etc.). Les problèmes actuels incluent la fragmentation des sources, le traitement manuel lent et le manque de traçabilité.

\subsection{Solution Proposée}
\textbf{Credit Platform} automatise ce processus via :
\begin{itemize}
    \item Centralisation multi-sources (Gmail, WhatsApp, Banque)
    \item Analyse NLP avancée pour extraire les données
    \item Classification par similarité vectorielle (IA)
\end{itemize}

\section{Architecture Système}

L'architecture modulaire repose sur 5 phases :

\begin{enumerate}
    \item \textbf{Ingestion} : Scripts Python pour récupérer les messages (Gmail API, Webhooks).
    \item \textbf{Traitement NLP} : Nettoyage et extraction d'entités (Montant, Client) avec spaCy.
    \item \textbf{Vectorisation (Cœur IA)} : Transformation des textes en vecteurs et recherche de similarité dans \textbf{Qdrant Cloud}.
    \item \textbf{Décision} : Moteur de règles métier basé sur l'intention détectée.
    \item \textbf{Dashboard} : Interface de visualisation et API FastAPI.
\end{enumerate}

\subsection{Intégration Qdrant (Focus)}
Nous n'utilisons pas de classification traditionnelle, mais une \textbf{recherche par similarité}.
Architecture technique :
\begin{itemize}
    \item \textbf{Base} : Qdrant Cloud
    \item \textbf{Collection} : \texttt{email\_memory}
    \item \textbf{Dimension} : 384 (all-MiniLM-L6-v2)
    \item \textbf{Distance} : Cosine
\end{itemize}

Chaque vecteur contient l'intention (\texttt{intent}) et le ton (\texttt{tone}) comme métadonnées, permettant d'identifier le but d'un nouveau message par comparaison avec l'historique.

\section{Pipeline de Données}

\subsection{Phase 1 & 2 : Ingestion et NLP}
Les messages sont récupérés, nettoyés (suppression HTML/signatures) et stockés dans MongoDB avec un statut \texttt{raw}. Le script NLP extrait ensuite les informations clés.

\subsection{Phase 3 : Vectorisation (Le Cœur)}
\begin{itemize}
    \item \textbf{Dataset Synthétique} : 80 emails réalistes générés pour servir de référence.
    \item \textbf{Processus} : Tout message entrant est vectorisé et comparé aux 80 exemples.
    \item \textbf{Résultat} : Les 5 plus proches voisins déterminent l'intention probable.
\end{itemize}

\subsection{Phase 4 & 5 : Décision et UI}
Selon l'intention (Ex: \texttt{CREDIT\_REQUEST}), un workflow spécifique est déclenché (réponse auto, création dossier).

\section{Timeline du Projet}

\subsection{Travaux Réalisés (Work Done)}
\begin{itemize}
    \item[$\checkmark$] \textbf{Infrastructure} : MongoDB, FastAPI, Structure du projet.
    \item[$\checkmark$] \textbf{Ingestion} : Scripts pour Gmail (`gmail\_fetch`), WhatsApp (`whatsapp\_fetch`) et Banque.
    \item[$\checkmark$] \textbf{NLP} : Nettoyage et extraction d'entités (`process\_messages.py`).
    \item[$\checkmark$] \textbf{Setup Qdrant} : Collection créée, connexion Cloud validée, Dataset synthétique généré.
\end{itemize}

\subsection{Travaux Futurs (Work to do)}
\begin{itemize}
    \item[$\square$] \textbf{Vectorisation Active} : Implémenter le script final de vectorisation (dépendance technique à résoudre).
    \item[$\square$] \textbf{Recherche Similarité} : Connecter le flux entrant à la recherche Qdrant en temps réel.
    \item[$\square$] \textbf{Moteur de Règles} : Câbler les actions métier (envoi email, etc.).
    \item[$\square$] \textbf{Dashboard} : Construire l'interface frontend finale.
\end{itemize}

\section{Conclusion}
L'architecture est en place. La connexion Qdrant est opérationnelle. La prochaine étape critique est la mise en production du pipeline de vectorisation pour activer l'intelligence du système.

\end{document}
